%% UNIT preface | IT authority lines 156--188
\pagenumbering{roman}
\section*{序言}
本书第二版的出版既已显得适宜，我认为，在保留全书基本构想的同时，最好大幅扩充其篇幅，使先前因篇幅过于狭窄而不得不压缩的许多理论部分，能够在新的论述中得到更为充分的发展。我还增添了若干新章：一方面用以论述那些虽已相当久远、却未能收入第一版的理论；另一方面用以介绍年代极近的研究，因为这些研究看来标志着无穷小几何在最近时期所取得的最重要进展。

本书的篇幅由此增加到原来的两倍，因而必须分为两卷；我力求使这种分卷大体上与一般理论和特殊理论的区分相对应。

本第一卷共有十四章。其中前十章在作了许多改动和增补之后，重现第一版相应各章的内容；我还从第一版录取如下说明。

\oldquote
“本课程所完全遵循的方法，发端于高斯著名的\emph{《曲面的一般研究》（Disquisitiones generales circa superficies curvas）》}；这一方法把无穷小几何看作对一个二次微分形式、或者同时对两个这样的形式所作的研究。因此，在论述双曲率曲线的第一章之后，读者将在第二章看到二次微分形式理论的基础；那里只作极为简要的陈述。
\closeoldquote

\oldquote
由这一理论产生的算法具有一项重大优点：它们在完全保留坐标线一般性的同时，使公式呈现出简洁优美、易于记忆的形式。本书后文始终使用这些算法，而本章正是其预备章节。

随后七章（第三至第九章）论述曲面的一般研究，也就是说，既研究曲面在空间中的实际形状所固有的性质，又研究那些在曲面弯曲时保持不变的性质（等距适用性理论）。第十章论述射线的二重无穷系统（直线汇）的一般理论；这一理论对于曲面的一般研究极为重要，无穷小几何的许多最卓越成果也由此而来。”
\closeoldquote

以下四章以更充分的展开和大量增补，重新呈现我曾在本书德文版\footnote{\emph{Vorlesungen über Differentialgeometrie}，M. Lukat 德译（Teubner，莱比锡，1899年）。}中给出的 \(n\) 维无穷小几何基本原理。第十一章先为任意 \(n\) 维空间的线度量和角度量规定基本约定，继而引入黎曼曲率的概念，并把通常理论中的公式推广到曲线和超曲面，尤其着重研究曲率线的性质。

然后我较为专门地研究具有常黎曼曲率的空间。首先在第十二章（相当于第一版的第十六章）讨论二维伪球几何，并对非欧几何作一简短说明。

在接下来的第十三章中，我一般地研究具有常曲率的 \(n\) 维空间，引入凯莱度量和魏尔斯特拉斯坐标系，并特别研究三维空间中运动的解析表示及其性质；最后讨论椭圆空间中重要的克利福德平行性概念，以及如何把曲线的弗雷内公式推广到椭圆几何和双曲几何。

最后，第十四章把这些一般理论应用于具有常曲率的 \(n\) 维空间中的超曲面，特别关注三维情形；在这一情形下，欧氏空间中已有的全部主要结果都得到推广。

我的同事 O. Niccoletti 教授\ednote{印本姓名拼作 ``Niccoletti''（双写 c）；继承的意大利文与英文转录曾无说明地写作 ``Nicoletti''。本简体中文版按印本恢复 ``Niccoletti''。}曾热心协助我完成校样的最后一次审订，谨在此向他致以最诚挚的谢意。

\begin{flushleft}
\emph{比萨，1902年4月。}
\end{flushleft}

\clearpage
\pagenumbering{arabic}
